\section{CREATE}
CREATE adopts a two-level architecture consisting of an outer reward-engineering agent and an underlying reinforcement-learning (RL) agent, as illustrated in Fig.~\ref{fig:framework}. The RL agent interacts with the environment and learns a policy under the current generated reward. The outer CREATE agent treats each completed policy-training and native-evaluation run as one reward-design round, calls diagnostic tools to observe training evidence, reflects on the current reward design, and edits the reward program.

\begin{figure*}[!t]
\centering
\includegraphics[width=0.98\textwidth]{create_framework.pdf}
\caption{Overview of CREATE. The outer reward-engineering agent observes training evidence produced by the underlying RL agent, reflects on the current reward design, and performs semantic-localized reward editing. Persistent lineage memory records intervention outcomes, while the best archive retains the reward program with the highest native-task performance.}
\label{fig:framework}
\end{figure*}

The resulting observe--reflect--act loop is written as
\begin{equation}
O_t=\Phi_{\mathcal T}(\mathcal D_t),\qquad
(H_t,A_t)=\Psi(R_t,O_t,M_t,B_t),\qquad
R_{t+1}=\mathcal U(R_t,A_t),
\label{eq:create_loop}
\end{equation}
where the training record $\mathcal D_t$ is transformed into observation $O_t$ through tool-augmented diagnosis. CREATE then combines $O_t$ with the current reward $R_t$, intervention memory $M_t$, and best-reward archive $B_t$ to produce structured reflection $H_t$, reward action $A_t$, and the next executable reward program $R_{t+1}$.

\subsection{Task Understanding and Reward Initialization}
Given a natural-language task description, the environment observation/action interface, and reward-code constraints, CREATE identifies the principal behavioral semantics and instantiates them as named reward components. The initial reward is
\begin{equation}
R_0(s,a,s')=\sum_{i=1}^{m_0}w_i^{(0)}r_i^{(0)}(s,a,s'),
\label{eq:reward_init}
\end{equation}
where each component $r_i^{(0)}$ represents a behavioral semantic, such as task progress, state stabilization, target events, or action cost, and $w_i^{(0)}$ determines its scale. Named components are logged independently while contributing to the total reward, enabling component-level diagnosis in later rounds. Generated code is checked for interface and execution validity before policy training; prompts and validation rules are provided in the supplementary material.

\subsection{Training-Evidence Observation, Reflection, and Action}
At iteration $t$, the underlying RL agent optimizes the generated reward, whereas the outer CREATE agent evaluates the resulting policy with the environment-native objective:
\begin{equation}
\pi_t=\arg\max_{\pi}\;\mathbb E_{\tau\sim\pi}\!\left[\sum_{k=0}^{H-1}\gamma^kR_t(s_k,a_k,s_{k+1})\right],\qquad
J_t=\frac{1}{N}\sum_{n=1}^{N}\sum_{k=0}^{H_n-1}r_{\rm env}^{(n,k)}.
\label{eq:train_eval}
\end{equation}
Thus, the generated reward shapes policy learning, while the native objective measures whether the learned behavior satisfies the task. The completed run produces $\mathcal D_t$, including episode outcomes, termination patterns, named-component activity and scale, and their changes over training. A diagnostic subagent queries these records through analysis tools and reports compact evidence to CREATE. Inactive components may indicate unreachable gates or sparse formulations; dominant components may overwhelm other objectives; and frequently active but weak components may require rescaling or reformulation. These patterns are diagnostic cues rather than fixed editing rules, leaving the reflection agent to infer interactions that are not captured by predefined heuristics.

CREATE then identifies one principal reward semantic and localizes the corresponding edit:
\begin{equation}
A_t=(q_t,\ell_t,\Delta_t),\qquad
\operatorname{supp}(\Delta_t)\subseteq\mathcal C(q_t).
\label{eq:semantic_edit}
\end{equation}
Here $q_t$ denotes the semantic under diagnosis, $\ell_t$ the edit level, and $\mathcal C(q_t)$ the components and parameters that implement that semantic. In most rounds, the action adds or removes one component, tunes its weight or gate, or changes its mathematical form. Related components may be adjusted jointly when they express the same semantic or require rebalancing, but independent design intents are not mixed within one ordinary revision. Level~1 calibrates coefficients, thresholds, and gates; Level~2 adds, removes, or refactors components while retaining the target semantic; Level~3 replans the reward structure after repeated local failure. This semantic-localized action space improves the traceability of reward interventions under stochastic and coupled RL training.

After the next run, CREATE records the previous observation, diagnosed semantic, edit, expected change, and observed outcome in $M_t$. The record supports comparison across rounds and discourages repeated unsuccessful interventions. In parallel, $B_t$ preserves the complete reward program with the highest native-task score, preventing later regressions from overwriting a previously discovered solution.

\subsection{Best-Reward Selection}
Because reward evolution is not necessarily monotonic, CREATE returns the best program observed within the reward-evaluation budget $K$:
\begin{equation}
t^*=\arg\max_{0\leq t<K}J_t,\qquad R^*=R_{t^*}.
\label{eq:best_reward}
\end{equation}
The selected reward is subsequently retrained and evaluated with independent policy and evaluation seeds to assess stability.